\section{Metodología}
\label{sec:metodologia}

\subsection{Tipo y Enfoque del Estudio}
\label{sec:met_tipo}

Esta investigación es de tipo \textbf{aplicada} y de enfoque \textbf{cuantitativo}, con
un diseño experimental comparativo. El paradigma subyacente es empírico-positivista y
el razonamiento sigue un esquema hipotético-deductivo: a partir del cuerpo teórico sobre
reconocimiento de emociones en voz se formulan hipótesis medibles, cuya corroboración
se sustenta en la cuantificación de métricas objetivas de rendimiento.

El alcance del estudio es \textbf{descriptivo-comparativo}: no solo se describe el
comportamiento individual de cada técnica, sino que se contrasta sistemáticamente el
rendimiento de tres métodos de \textit{baseline} cruzados con cuatro algoritmos de
detección, lo que genera un diseño factorial $3 \times 4$ sobre los mismos datos de
evaluación.

\subsection{Diseño Experimental}
\label{sec:met_diseno}

\subsubsection{Variables independientes}

\begin{itemize}
    \item \textbf{Método de \textit{baseline} individual} $\in$ \{inicio (primeros
          2--5~min de llamada), mediana global, ventana deslizante adaptativa\}.
    \item \textbf{Algoritmo de detección de anomalías} $\in$ \{Z-score estadístico,
          Isolation Forest, One-Class SVM, Random Forest\}.
\end{itemize}

\subsubsection{Variables dependientes}

Precision, Recall, F1-score y Accuracy, calculadas para la tarea de detección de picos
emocionales (clase positiva: pico; clase negativa: segmento neutro) y para la
clasificación de valencia emocional (positiva/negativa).

\subsubsection{Control experimental}

Se emplea validación cruzada estratificada con $k = 5$ particiones para todas las
combinaciones del diseño factorial, garantizando que la distribución de clases se
mantenga en cada fold. Adicionalmente, los resultados se contrastan con anotaciones
de expertos humanos como referencia externa independiente.

\subsubsection{Hipótesis de trabajo}

La combinación \textit{baseline} adaptativa con ventana deslizante + Isolation Forest
produce el mayor F1-score en la detección de picos emocionales sobre el corpus de
llamadas reales, superando estadísticamente a los métodos con \textit{baseline} fija
($\alpha = 0{,}05$, prueba de Wilcoxon de rangos con signo \cite{wilcoxon1945}).

\subsection{Población y Muestra}
\label{sec:met_poblacion}

\subsubsection{Datasets}

\textbf{Entrenamiento y calibración:} Se utiliza RAVDESS (\textit{Ryerson Audio-Visual
Database of Emotional Speech and Song}), corpus actuado con 24 actores profesionales
(12 hombres, 12 mujeres) que cubre 8 categorías emocionales. Aunque el habla actuada
difiere de la espontánea, RAVDESS ofrece etiquetado de referencia limpio para calibrar
los algoritmos antes de exponerlos al corpus real \cite{abdelhamid2022}.

\textbf{Evaluación y validación:} Corpus propio de llamadas reales de call center en
español (Colombia). Se establece como objetivo mínimo un corpus de 30 llamadas (no
menos de 5 minutos cada una), lo que representa aproximadamente 150 minutos de audio
anotado. Los criterios de inclusión son: (a) calidad de audio aceptable (SNR $\geq
15$~dB), (b) duración mínima de 5~minutos, (c) presencia de ambos interlocutores
(agente y cliente), y (d) disponibilidad de consentimiento para uso en investigación.
Los criterios de exclusión son: (a) ruido excesivo o degradación del códec, (b)
interrupciones técnicas superiores a 30~segundos, y (c) llamadas en las que uno de los
interlocutores no habla en español estándar colombiano.

\subsubsection{Proceso de anotación}

La anotación será realizada por mínimo dos expertos independientes (psicólogos o
fonoaudiólogos con experiencia en análisis emocional). Cada experto escuchará la
llamada completa y marcará los \textit{timestamps} de inicio y fin de cada pico
emocional, junto con la etiqueta de valencia (positiva o negativa). El acuerdo
inter-anotador se cuantificará con el coeficiente kappa de Cohen \cite{cohen1960kappa};
se aceptarán como anotaciones definitivas únicamente los segmentos con $\kappa \geq
0{,}7$, umbral considerado como acuerdo sustancial en la literatura \cite{cohen1960kappa}.
Los desacuerdos se resolverán mediante consenso supervisado.

\subsection{Herramientas e Instrumentos}
\label{sec:met_herramientas}

\subsubsection{Hardware}

Procesador Intel Core i7-4770 (Dell Optiplex 9020), 16~GB RAM, sistema operativo
Linux. Las restricciones de cómputo justifican la elección de métodos \textit{feature-based}
clásicos frente a modelos de aprendizaje profundo de gran escala.

\subsubsection{Software}

\begin{itemize}
    \item \textbf{Python 3.8+}: lenguaje principal de implementación.
    \item \textbf{Librosa}: extracción de F0, intensidad, MFCCs, tasa de habla.
    \item \textbf{Parselmouth}: bindings Python de Praat para el cálculo de jitter
          y shimmer.
    \item \textbf{OpenSMILE + eGeMAPS}: conjunto de 88 parámetros acústicos
          minimizados validados para investigación en voz y computación afectiva
          \cite{eyben2016gemaps}.
    \item \textbf{pyannote.audio}: diarización de hablantes para separar los canales
          de agente y cliente.
    \item \textbf{scikit-learn}: implementación de Isolation Forest, One-Class SVM,
          Random Forest, validación cruzada estratificada y métricas de evaluación.
    \item \textbf{NumPy, Pandas, Matplotlib}: manipulación de datos y visualización
          de resultados.
\end{itemize}

El control de versiones se gestiona con Git; el repositorio público del proyecto está
disponible en \url{https://github.com/PirrY/Desviaci-nesEmocionalesAudio}.

\subsection{Procedimiento}
\label{sec:met_procedimiento}

El sistema sigue un pipeline secuencial de seis etapas, alineadas con las seis fases
del cronograma de trabajo (véase la Tabla~\ref{tab:cronograma}). La
Figura~\ref{fig:pipeline} ilustra el flujo de datos.

\begin{figure*}[!t]
\centering
\begin{tikzpicture}[
    node distance=0.45cm and 0.35cm,
    box/.style={
        rectangle, draw=black, rounded corners=3pt,
        minimum height=1.0cm, text width=1.95cm,
        align=center, font=\scriptsize, fill=white
    },
    arrow/.style={-{Stealth[length=4pt]}, thick}
]
  \node[box] (audio)  {Audio\\crudo};
  \node[box, right=of audio]  (prep)  {Prepro-\\cesamiento};
  \node[box, right=of prep]   (feat)  {Extracción\\de \textit{features}};
  \node[box, right=of feat]   (base)  {Cálculo de\\\textit{baseline}};
  \node[box, right=of base]   (det)   {Detector\\de picos};
  \node[box, right=of det]    (val)   {Clasificador\\de valencia};
  \node[box, right=of val]    (out)   {Salida\\(timestamps)};

  \draw[arrow] (audio) -- (prep);
  \draw[arrow] (prep)  -- (feat);
  \draw[arrow] (feat)  -- (base);
  \draw[arrow] (base)  -- (det);
  \draw[arrow] (det)   -- (val);
  \draw[arrow] (val)   -- (out);
\end{tikzpicture}
\caption{Pipeline del sistema de detección de picos emocionales. El flujo va desde el
audio crudo hasta la generación de \textit{timestamps} etiquetados con su valencia.}
\label{fig:pipeline}
\end{figure*}

\textbf{Fase 1 — Configuración y extracción de \textit{features} (sem.\ 1--3):}
Instalación del entorno, descarga de RAVDESS, implementación del módulo de extracción
acústica (F0, intensidad, jitter, shimmer, MFCCs, tasa de habla).

\textbf{Fase 2 — \textit{Baseline} emocional y preprocesamiento (sem.\ 4--5):}
Implementación de los tres métodos de \textit{baseline} individual. Diarización de
hablantes con pyannote.audio para separar los canales de agente y cliente.

\textbf{Fase 3 — Detección de picos emocionales (sem.\ 6--8):}
Entrenamiento y comparación de los cuatro algoritmos de detección (Z-score, Isolation
Forest, One-Class SVM, Random Forest) sobre el corpus de calibración.

\textbf{Fase 4 — Clasificación de valencia emocional (sem.\ 9--10):}
Desarrollo del clasificador binario de valencia. Optimización de hiperparámetros
mediante búsqueda en cuadrícula (\textit{grid search}).

\textbf{Fase 5 — Validación y experimentación (sem.\ 11--13):}
Ejecución del sistema completo sobre las llamadas reales anotadas. Análisis comparativo
del diseño factorial $3 \times 4$.

\textbf{Fase 6 — Documentación e integración final (sem.\ 14--15):}
Consolidación del código, documentación y preparación del informe final.

\subsection{Técnicas de Análisis}
\label{sec:met_analisis}

\subsubsection{Métricas de evaluación}

Las métricas primarias son Precision, Recall, F1-score y Accuracy, definidas a partir
de la matriz de confusión binaria (pico vs.\ no-pico; positivo vs.\ negativo). Se
generarán curvas ROC y se calculará el área bajo la curva (AUC) para el detector
binario pico/no-pico, permitiendo evaluar el \textit{trade-off} entre la tasa de
verdaderos positivos y la de falsos positivos.

\subsubsection{Importancia de características}

Se analizará la importancia de las \textit{features} mediante el atributo
\texttt{feature\_importances\_} del modelo Random Forest, identificando cuáles
características acústicas son más discriminativas para la detección de picos y la
clasificación de valencia.

\subsubsection{Pruebas de significancia}

La comparación estadística entre algoritmos se realizará con la prueba de Wilcoxon de
rangos con signo \cite{wilcoxon1945} aplicada sobre los F1-scores obtenidos en los
cinco folds de validación cruzada. Se utilizará un nivel de significancia $\alpha =
0{,}05$. Esta prueba no paramétrica es adecuada dado que no se puede asumir
normalidad en la distribución de las métricas con $k = 5$ \cite{wilcoxon1945}.

\subsection{Consideraciones Éticas}
\label{sec:met_etica}

\begin{itemize}
    \item \textbf{Anonimización:} antes de su uso en investigación, las grabaciones
          serán procesadas para eliminar nombres, números de identificación, datos
          bancarios y cualquier otro dato personal identificable.
    \item \textbf{Consentimiento informado:} la empresa proveedora de las llamadas
          otorgará consentimiento escrito para el uso académico de las grabaciones,
          en cumplimiento de los principios de finalidad y libertad establecidos en
          la Ley 1581 de 2012 de la República de Colombia \cite{ley1581colombia}.
    \item \textbf{Almacenamiento seguro:} los archivos de audio originales se
          almacenarán con cifrado en reposo y acceso restringido.
    \item \textbf{Habeas Data:} el tratamiento de datos de voz se enmarca en la Ley
          1581 de 2012 \cite{ley1581colombia}, que regula la protección de datos
          personales en Colombia.
    \item \textbf{Conflicto de interés:} el autor declara no tener ningún conflicto
          de interés con las organizaciones proveedoras de datos.
\end{itemize}

\subsection{Plan de Validación}
\label{sec:met_validacion}

\subsubsection{Protocolo de anotación}

Los expertos recibirán una guía de anotación estandarizada que incluye: (a) definición
operacional de \textit{pico emocional} (incremento sostenido de al menos dos
características acústicas por encima de $1{,}5\sigma$ respecto al \textit{baseline}
individual durante un mínimo de 5~segundos); (b) escala de valencia binaria
(positiva: alegría, entusiasmo, satisfacción; negativa: ira, frustración, ansiedad);
y (c) procedimiento para marcar \textit{timestamps} en Praat o Audacity.

\subsubsection{Comparación sistema--expertos}

Las detecciones automáticas se compararán con las anotaciones humanas usando una
tolerancia temporal de $\pm 3$~segundos en los bordes de los segmentos. Un segmento
detectado por el sistema se considerará verdadero positivo si su centro coincide con
el de una anotación de experto dentro de dicha tolerancia.

\subsubsection{Criterios de aceptación}

El sistema se considerará exitoso si alcanza, sobre el corpus de llamadas reales
anotadas, los umbrales definidos en la Sección~\ref{sec:resultados_esperados}:
Precision $\geq 75$\,\%, Recall $\geq 70$\,\%, F1-score $\geq 72$\,\%.
